\section*{Chapitre \ref{ch:ml:representation-learning} --- Apprentissage de représentations}

\begin{solution}{exo:ml:representation-learning:1}
Parmi $2m - 1 = 511$ candidats~; avec des codes non informatifs, chaque candidat est également probable et la perte vaut
$\ln 511 = 6.24$.
\end{solution}

\begin{solution}{exo:ml:representation-learning:2}
$300\times0.5 = 150$ secondes, 2.5 minutes~; $4\,444\times0.5 = 2\,222$ secondes, 37.0 minutes.
\end{solution}

\begin{solution}{exo:ml:representation-learning:3}
$400\times64 + 64 + 64\times16 + 16 = 26\,704$ paramètres~; une \omterm{def:ml:representation-learning:probe}{sonde linéaire} sur le code en a $16 + 1 = 17$.
\end{solution}

\begin{solution}{exo:ml:representation-learning:4}
Un objectif de reconstruction conserve les directions de plus forte variance (la \cref{prop:ml:represent:pca} pour le cas linéaire),
et, sur un carnet d'ordres, ce sont les tailles des niveaux profonds~; le déséquilibre au meilleur niveau, qui prédit les secondes suivantes, ne représente qu'une faible
part de la variance et il est en partie perdu dans un code à 16 dimensions. La fenêtre brute le conserve, et une régression ridge peut le retrouver.
\end{solution}

\begin{solution}{exo:ml:representation-learning:5}
Une régression ridge sur 400 entrées brutes a beaucoup de paramètres à estimer~; avec peu d'étiquettes, ses estimations sont bruitées et un code à 16 dimensions
qui pointe déjà vers les bonnes directions l'emporte. Avec 4\,444 étiquettes, la régression ridge estime elle-même les directions utiles,
et la compression du code commence à coûter plus qu'elle ne rapporte.
\end{solution}

\begin{solution}{exo:ml:representation-learning:6}
Le \omterm{def:ml:representation-learning:ssl}{pré-entraînement} sur les mois de test laisse l'encodeur ajuster leurs états particuliers~: une fuite transductive qui
n'exige aucune étiquette (le décalage de distribution sur la période de test est absorbé dans la représentation). Pré-entraîner uniquement sur des données antérieures à la
période de test, comme pour tout composant ajusté.
\end{solution}

\begin{solution}{exo:ml:representation-learning:7}
\texttt{compare(1000, ahead=16)}~: la sonde prédictive tombe à 0.204, contre 0.285 à deux secondes, au niveau du réseau
entraîné depuis zéro. À huit secondes, le carnet est en grande partie imprévisible, si bien que la cible du prétexte est surtout du bruit et que
l'encodeur en apprend moins.
\end{solution}

\begin{solution}{exo:ml:representation-learning:8}
Pour la vue $a$ de partenaire $p(a)$, le terme de perte vaut $-\ln\bigl(e^{s_{a,p(a)}}/\sum_{b\ne a}e^{s_{ab}}\bigr)$~:
l'entropie croisée d'un softmax sur les $2m - 1$ autres vues, le partenaire étant la classe correcte. Si toutes les similarités sont
égales, chaque terme vaut $-\ln(1/(2m-1)) = \ln(2m - 1)$, et leur moyenne aussi.
\end{solution}

\begin{solution}{pb:ml:representation-learning:1}
\begin{enumerate}
\item Reconstruire une fenêtre masquée et bruitée~; faire correspondre deux vues augmentées d'une même fenêtre à d'autres fenêtres~; prévoir l'évolution
du carnet deux secondes après la fenêtre.
\item Le prétexte de prévision~: sa cible est le futur proche du flux lui-même, du même type d'objet que celui de la tâche.
\item Que son code conserve les directions de plus forte variance, qui ne sont pas forcément prédictives.
\item Que le bruit gaussien et le masquage aléatoire des entrées ne modifient pas ce qui compte dans une fenêtre.
\item 0.22, 0.12 et 0.18.
\item 0.29, 0.24 et 0.20.
\item Moins bon que la fenêtre brute à 1\,000 et à 4\,444~; à 300, l'\omterm{def:ml:cross-sectional-deep-models:ae}{autoencodeur} donne 0.15, l'encodeur contrastif 0.05.
\item Dix variables ajustées sur deux minutes et demie d'un flux surtout plat sont du bruit~; avec davantage de données, elles portent le
savoir de microstructure que les réseaux doivent apprendre.
\item Régression ridge sur la fenêtre brute 0.45, sonde prédictive 0.40, \omterm{def:ml:representation-learning:ssl}{ajustement fin} 0.46, à partir de zéro 0.38, l'\omterm{def:ml:cross-sectional-deep-models:ae}{autoencodeur} et
les sondes contrastives 0.31, les variables faites à la main 0.55.
\item L'\omterm{def:ml:representation-learning:ssl}{ajustement fin} adapte le code de l'encodeur à la tâche dès qu'il existe assez d'étiquettes pour le faire sans risque~; un code figé conserve
seulement ce que le prétexte a retenu.
\item $0.46 - 0.38 = 0.08$ d'IC.
\item Les courbes tendent à converger~: avec beaucoup d'étiquettes, l'avantage du \omterm{def:ml:representation-learning:ssl}{pré-entraînement} s'estompe et la régression ridge sur les entrées brutes ou faire à la main rattrape
son retard.
\item \emph{Résultat à retenir}~: en dessous d'environ 1\,000 fenêtres étiquetées (huit minutes), la sonde prédictive bat tout modèle à entrées brutes
entraîné sur les seules étiquettes~: 0.22 contre 0.18 à 300 étiquettes en partant de zéro, 0.29 contre 0.20 à 1\,000~; avec 4\,444, seul
l'encodeur ajusté finement fait jeu égal avec la régression ridge.
\item Lorsque les étiquettes sont rares et les données non étiquetées abondantes, et lorsque personne n'a encore formalisé les variables (un nouveau marché,
un nouveau type de données).
\item Ce qui ne doit pas modifier une prévision~: un faible bruit de mesure, des décalages d'un instantané, un recalibrage des niveaux profonds~; il faut les tester
avec une sonde, et non à l'intuition.
\item La période de test et tout ce qui la suit.
\item Regrouper les sessions de nombreux instruments, avec une normalisation par instrument, et sonder chacun d'eux.
\item Quelle part de la tâche est accessible linéairement dans la représentation, indépendamment de la manière dont un modèle aval est
entraîné.
\item C'est un très grand modèle auto-supervisé de prédiction du texte, utilisé au moyen de sondes, de prompts ou d'un \omterm{def:ml:representation-learning:ssl}{ajustement fin}.
\item Être proche de la tâche, de sorte que ce qu'il conserve soit ce dont la tâche a besoin.
\end{enumerate}
\end{solution}

\begin{solution}{iq:ml:representation-learning:1}
Une \omterm{def:ml:representation-learning:probe}{sonde linéaire} ajuste un modèle linéaire sur des codes figés~: peu coûteuse, elle mesure ce que contient la représentation, et reste sûre avec peu d'étiquettes.
L'\omterm{def:ml:representation-learning:ssl}{ajustement fin} entraîne aussi l'encodeur~: il peut adapter la représentation, demande davantage d'étiquettes et de précautions (faible \omterm{def:ml:trees-and-boosting:boosting}{taux d'apprentissage}, \omterm{def:ml:trees-and-boosting:early}{arrêt
précoce}). Commencer par la sonde~; ajuster finement lorsqu'il y a assez d'étiquettes.
\iqlookfor{le compromis entre adaptation et \omterm{def:ml:why-financial-machine-learning-is-different:generalisation}{surapprentissage}.}
\end{solution}

\begin{solution}{iq:ml:representation-learning:2}
Entraîner des codes de sorte que deux vues d'un même exemple soient proches et que les vues d'exemples différents soient éloignées (InfoNCE). Pour les prix
intrajournaliers~: un faible bruit, de petits décalages temporels, un masquage de certaines variables, un recalibrage des volumes~; jamais de transformations qui changent
le sens ou le résultat futur (inverser le temps, échanger achats et ventes).
\iqlookfor{l'objectif et les augmentations qui respectent ce qui compte.}
\end{solution}

\begin{solution}{iq:ml:representation-learning:3}
Pré-entraîner un encodeur sur les dix ans avec un prétexte auto-supervisé proche de la tâche (prévoir les prochaines secondes ou
jours du flux), puis sonder et éventuellement ajuster finement sur les trois mois~; comparer sur des données étiquetées mises de côté avec un modèle
entraîné à partir de zéro.
\iqlookfor{un prétexte aligné sur la tâche et une comparaison honnête.}
\end{solution}

\begin{solution}{iq:ml:representation-learning:4}
Il conserve ce qui varie le plus dans l'entrée (ACP pour un \omterm{def:ml:cross-sectional-deep-models:ae}{autoencodeur} linéaire), et la partie prédictive de l'entrée peut être une
direction de faible variance que le code abandonne.
\iqlookfor{la variance n'est pas la pertinence.}
\end{solution}

\begin{solution}{iq:ml:representation-learning:5}
Des \omterm{def:ml:representation-learning:probe}{sondes linéaires} sur données mises de côté pour plusieurs tâches, et la courbe d'apprentissage sur données étiquetées, face à une référence à entrées brutes.
\iqlookfor{sondes et courbes d'apprentissage.}
\end{solution}

\begin{solution}{iq:ml:representation-learning:6}
La reconstruction $BAX$ est de rang au plus $k$~; par Eckart--Young, la meilleure approximation de rang-$k$ de $X$ au sens de Frobenius est la
décomposition en valeurs singulières tronquée, c'est-à-dire la projection sur les $k$ premiers vecteurs singuliers à droite, soit les directions principales~; choisir $A$ comme
les coordonnées de la projection et $B$ comme sa base permettent de l'atteindre.
\iqlookfor{l'argument de rang et Eckart--Young.}
\end{solution}
